\section{Giới thiệu}
Hệ thống khử nhiễu tiếng nói đơn kênh (Single-channel Speech Enhancement) luôn là một bài toán kinh điển nhưng đầy thách thức trong lĩnh vực xử lý tín hiệu. Mặc dù các kỹ thuật truyền thống như trừ phổ (spectral subtraction) \cite{oppenheim1978evaluation} hay lọc hậu kỳ (post-filtering) \cite{chen1995adaptive} đã được ứng dụng rộng rãi, chúng thường để lại những vệt nhiễu cơ học (musical noise) do suy giảm năng lượng không đồng đều giữa các dải tần. Sự bùng nổ của Deep Learning mở ra cơ hội giải quyết bài toán này với độ phân giải cao hơn. Tuy nhiên, các kiến trúc mô hình học sâu truyền thống thường yêu cầu mô hình hóa hàng ngàn điểm tần số đồ sộ, dẫn đến khối lượng tính toán "khổng lồ", không phù hợp với các hệ thống nhúng với ràng buộc thời gian thực (real-time constraints) \cite{valin2018hybrid}. 

Khác với cách tiếp cận End-to-End, mô hình RNNoise nguyên bản của J.-M. Valin \cite{valin2018hybrid} giải quyết rào cản tính toán bằng cách kết hợp sức mạnh phân loại của Deep Learning với độ tinh giản của xử lý tín hiệu số (DSP). Thay vì tái tạo lại toàn bộ dạng sóng âm, hệ thống học sâu chỉ tập trung dự báo các hệ số khuếch đại (ideals gains) trên các dải băng tần thính giác (Bark scale), nhường phần tinh chỉnh tín hiệu giữa các hài âm (harmonics) cho hệ thống lọc lược truyền thống (pitch comb filtering). Mô hình thực tế triển khai trong báo cáo này đã được nâng cấp với 32 dải Bark (thay vì 22 như \cite{valin2018hybrid}) và đưa thêm các khối chập (1D-CNN) nhằm nâng cao sức mạnh nhận diện trên không gian \cite{valin2024noise}. 

Dù mô hình này đã rất tinh gọn gọn nhẹ (~1.3 triệu tham số dạng INT8 đối với cấu hình thực tế), việc triển khai suy luận với tốc độ 10ms/frame (đối với mẫu 48kHz) trên phần cứng ESP32-S3 đặt ra giới hạn nghiêm ngặt không chỉ về số phép toán dấu phẩy động (FLOP/s) mà còn về cường độ I/O của bộ nhớ (Memory Bandwidth).

Bài báo cáo này được cấu trúc lại như sau: Phần 2 mô tả tường minh luồng tín hiệu lai và mô hình học sâu. Phần 3 tập trung phân tích các phương pháp tối ưu hóa cực đoan (SIMD PIE, IRAM, Quantization) đối với ESP32-S3, đồng thời chỉ ra rào cản vật lý của bộ đệm (PSRAM Bottleneck). Phần 4 đề xuất thiết kế kiến trúc phân tán Edge-to-Server để triệt tiêu giới hạn phần cứng, và Phần 5 kết luận lại ưu, nhược điểm của giải pháp đạt được, đi từ phần cứng nhúng đến vi kiến trúc cấp hệ thống.
